% from ega2-1-fr.tex
\oldpage[II]{24}
l'anneau $(S^{(d)})_f$, qui est donc \emph{isomorphe à
$S_{(f)}[T,T^{-1}]=S_{(f)}\otimes_{\mathbf{Z}}\mathbf{Z}[T,T^{-1}]$}
(où $T$ est une indéterminée). En effet, si on a une relation
$\sum_{h=-a}^{b}z_h(f/1)^h=0$ avec $z_h=x_h/f^m$, où tous les $x_h$
appartiennent à $S_{md}$, cette relation équivaut par définition à
l'existence d'un $k\geq a$ tel que
$\sum_{h=-a}^{b}f^{h+k}x_h=0$, et comme les degrés des termes de cette somme
sont distincts, on a $f^{h+k}x_h=0$ pour tout $h$, d'où $z_h=0$ pour tout
$h$.

Si $M$ est un $S$-module gradué, $M'=M_f$ est un $S_f$-module gradué,
$M'_n$ étant l'ensemble des $z/f^k$ avec $z\in M_{n+kd}$ ($k\geq 0$)~; on
désigne par $M_{(f)}$ l'ensemble des éléments homogènes de degré $0$ de
$M'$~; il est immédiat que $M_{(f)}$ est un $S_{(f)}$-module et que l'on a
$(M^{(d)})_f=M_{(f)}\otimes_{S_{(f)}}(S^{(d)})_f$.
\end{env}

\begin{lemma}[2.2.2]
\label{II.2.2.2-fr}
Soient $d$, $e$ deux entiers $>0$, $f\in S_d$, $g\in S_e$. Il existe un
isomorphisme canonique d'anneaux
\[
  S_{(fg)}\xrightarrow{\sim}(S_{(f)})_{g^d/f^e}~;
\]
si on identifie canoniquement ces deux anneaux, il existe un isomorphisme
canonique de modules
\[
  M_{(fg)}\xrightarrow{\sim}(M_{(f)})_{g^d/f^e}.
\]
\end{lemma}

\begin{proof}
En effet, $fg$ divise $f^e g^d$, et ce dernier élément divise $(fg)^{de}$,
donc les anneaux gradués $S_{fg}$ et $S_{f^e g^d}$ s'identifient
canoniquement~; d'autre part, $S_{f^e g^d}$ s'identifie aussi à
$(S_{f^e})_{g^d/1}$ (\hyperref[0.1.4.6-fr]{0, I.4.6}), et comme $f^e/1$ est
inversible dans $S_{f^e}$, $S_{f^e g^d}$ s'identifie aussi à
$(S_{f^e})_{g^d/f^e}$. Or l'élément $g^d/f^e$ est de degré $0$ dans
$S_{f^e}$~; on en conclut aussitôt que le sous-anneau de
$(S_{f^e})_{g^d/f^e}$ formé des éléments de degré $0$ est
$(S_{(f^e)})_{g^d/f^e}$, et comme on a évidemment
$S_{(f^e)}=S_{(f)}$, cela démontre la première partie de la proposition~;
la seconde s'établit de même.
\end{proof}

\begin{env}[2.2.3]
\label{II.2.2.3-fr}
Avec les hypothèses de (\hyperref[II.2.2.2-fr]{2.2.2}), il est clair que
l'homomorphisme canonique $S_f\to S_{fg}$
(\hyperref[0.1.4.1-fr]{0, I.4.1}), qui transforme $x/f^k$ en
$g^k x/(fg)^k$ est de degré $0$, donc donne par restriction un
\emph{homomorphisme canonique $S_{(f)}\to S_{(fg)}$}, tel que le diagramme
\[
  \xymatrix{
    & S_{(f)}\ar[dl]\ar[dr]\\
    S_{(fg)}\ar[rr]^-{\sim} && (S_{(f)})_{g^d/f^e}
  }
\]
soit commutatif. On définit de même un homomorphisme canonique
$M_{(f)}\to M_{(fg)}$.
\end{env}

\begin{lemma}[2.2.4]
\label{II.2.2.4-fr}
Si $f$, $g$ sont deux éléments homogènes de $S_+$, l'anneau $S_{(fg)}$ est
engendré par la réunion des images canoniques de $S_{(f)}$ et $S_{(g)}$.
\end{lemma}

\begin{proof}
En vertu de (\hyperref[II.2.2.2-fr]{2.2.2}), il suffit de voir que
$1/(g^d/f^e)=f^{d+e}/(fg)^d$ appartient à l'image canonique de $S_{(g)}$
dans $S_{(fg)}$, ce qui est évident par définition.
\end{proof}

\begin{proposition}[2.2.5]
\label{II.2.2.5-fr}
Soit $d$ un entier $>0$ et soit $f\in S_d$. Alors il existe un isomorphisme
canonique d'anneaux
$S_{(f)}\xrightarrow{\sim}S^{(d)}/(f-1)S^{(d)}$~; si on identifie ces deux
anneaux par cet isomorphisme, il existe un isomorphisme canonique de modules
$M_{(f)}\xrightarrow{\sim}M^{(d)}/(f-1)M^{(d)}$.
\end{proposition}

\begin{proof}
Le premier de ces isomorphismes se définit en faisant correspondre à
$x/f^n$, où $x\in S_{nd}$, l'élément $\overline{x}$, classe de $x$ mod.
$(f-1)S^{(d)}$~; cette application est bien définie, car on a la congruence
$f^h x\equiv x\pmod{(f-1)S^{(d)}}$ pour tout $x\in S^{(d)}$, donc si
$f^h x=0$ pour un $h>0$,
